\paragraph{Edge of a toric block on a product of two fuzzy spheres (conditional; bounded).}
Take the six Hermitian \(N^2\times N^2\) matrices \(X^I=(A_1{}^a{}_bJ^b\otimes1,\;1\otimes A_2{}^a{}_bJ^b)\), with \(A_1,A_2\) real and invertible, \(J^b\) the \(N\)-dimensional irreducible representation of \(\mathfrak{su}(2)\) and the positive contraction \(\delta_{IJ}\); this is a compact four-dimensional fuzzy factor. Let \(|i_1,i_2)\), \(0\le i_f\le N-1\), be the joint eigenvectors of \(J^3\otimes1\) and \(1\otimes J^3\), with eigenvalues \(\tfrac12(N-1)-i_f\), and let \(\Theta\) be the projector on a set \(\Omega\) of sites \((i_1,i_2)\), called here a toric block; the spectral projectors of \(a\,J^3\otimes1+b\,1\otimes J^3\) are of this form. In the Poisson limit the edge of such a block is three-dimensional, the dimension of \(\partial M\) in Definition~\ref{def:HBEV}, whereas the cap of the preceding paragraph has a one-dimensional edge.

\emph{Edge operator.} Suppose that no line of sites in either lattice direction meets \(\Omega\) in the pattern inside, outside, inside on three consecutive sites. Then \(E\equiv\sum_I\Theta X^I\bar\Theta X^I\Theta\) is diagonal in \(|i_1,i_2)\), and its entry at a site of \(\Omega\) is the sum, over the nearest-neighbour links from that site to sites outside \(\Omega\), of \(\kappa=k_f\,i(N-i)\) for a link between \(i-1\) and \(i\) in factor \(f\), with \(k_f=\tfrac14(\sigma_{f,1}^2+\sigma_{f,2}^2)\) and \(\sigma_{f,1,2}\) the singular values of \(A_f\) on the plane orthogonal to the third axis (exact). So the rank of the map \(\mathcal E=(\Theta X^I\bar\Theta)_I\) is the number of edge sites, the sites of \(\Omega\) with a neighbour outside; \(s^2=\operatorname{Tr}E\) is the squared norm of the preceding paragraph, and \(T\equiv\operatorname{Tr}E^{1/2}\) is the trace norm of \(\mathcal E\). If the pattern occurs in factor \(f\), \(E\) has an entry between the two inside sites proportional to \(\sum_a(A_f{}^a{}_1-\mathrm iA_f{}^a{}_2)^2\), which vanishes only for \(\sigma_{f,1}=\sigma_{f,2}\) (exact). For the cap times the whole second sphere, \(\Omega=\{i_1<M_1\}\), the rank is \(N\) and \(E=s_1^2\) on each of the \(N\) edge sites, with \(s_1^2=k_1M_1(N-M_1)\), so
\[
T=N\,s_1,\qquad s^2=N\,s_1^2,
\]
both independent of \(A_2\) (exact). For the product block \(\Omega=\{i_1<M_1,\ i_2<M_2\}\) with \(M_f<N\) and \(s_f^2=k_fM_f(N-M_f)\), the rank is \(M_1+M_2-1\), \(s^2=M_2s_1^2+M_1s_2^2\) and \(T=(M_2-1)s_1+(M_1-1)s_2+(s_1^2+s_2^2)^{1/2}\) (exact).

\emph{Measure.} Let \(x^I\) be the symbols, \(\theta^{\alpha\beta}\) the Poisson bivector, \(\rho=1/|\operatorname{Pf}\theta|\), and let \(g_{\alpha\beta}\) and \(K^{\mu\nu}\) be the induced tensor and the operator tensor defined as in Definition~\ref{def:matrix-coords-ext} with \(\delta_{IJ}\) in place of \(\eta_{IJ}\). For a function \(f\), \(K^{\mu\nu}\partial_\mu f\partial_\nu f=\sum_I\{f,x^I\}^2\), and
\[
P_K\equiv(2\pi)^{-2}\int\mathrm d^4y\,\rho\,\delta(f)\,\bigl(K^{\mu\nu}\partial_\mu f\,\partial_\nu f\bigr)^{1/2}
\]
is \((2\pi)^{-2}\) times the three-volume of \(\{f=0\}\) in the tensor \(\tilde\gamma_{\mu\nu}=(\sqrt g/\rho)^{2/3}(K^{-1})_{\mu\nu}\), with \(\sqrt g=(\det g_{\alpha\beta})^{1/2}\) (exact, a pointwise identity; the normalisation is one site per volume \((2\pi)^2\) of \(\rho\,\mathrm d^4y\)). For the cap times the whole second sphere, with \(M_1(N-M_1)\to(\tfrac12N\sin\theta)^2\) as in the preceding paragraph, \(\sqrt2\,T/P_K\) is the ratio of the root mean square to the mean of the speed \(|\partial_{\phi_1}x|\) over the angle \(\phi_1\) of the flow that generates the cut; it is at least \(1\), with equality only for \(\sigma_{1,1}=\sigma_{1,2}\). Under \(A_1\to\lambda A_1\) and under \(A_2\to\lambda A_2\), \(T\) and \(P_K\) scale as \(\lambda\) and \(\lambda^0\), the three-volume of the edge in \(g\) as \(\lambda\) and \(\lambda^2\), and that in \(K^{-1}\) as \(\lambda^{-1}\) and \(\lambda^{-2}\) (exact). So \(T\) weighs the directions of the edge that are transverse to that flow by the number of sites, that is by the symplectic volume, and only the direction of the flow by a length in the induced tensor; it is the three-volume of the edge neither in \(g\) nor in \(K^{-1}\). For a half-plane block \(a\,i_1+b\,i_2<c\) with \(a\ge b\ge0\) on the round product, \(A_f=\nu_f1\), the edge sites are the last sites of the lines of fixed \(i_2\), and such a site has a second outside neighbour where the boundary steps. Per unit Euclidean length of the boundary line in the plane of \((i_1,i_2)\), with unit normal \((n_1,n_2)\) and \(e_f=k_fi_f(N-i_f)\), the number of edge sites is therefore \(n_1\), the density of \(s^2\) is \(n_1e_1+n_2e_2\), and the density of \(T\) is
\[
(n_1-n_2)\,e_1^{1/2}+n_2\,(e_1+e_2)^{1/2}\;\ge\;\bigl(n_1^2e_1+n_2^2e_2\bigr)^{1/2},
\]
the right-hand side being the density of \(P_K/\sqrt2\); equality holds only for \(n_2=0\) and for \(n_1=n_2\). At \(N=3200\) the lattice sums agree with the integrals of these densities to \(10^{-4}\), and \(\sqrt2\,T/P_K\) is \(1.113\) and \(1.096\) for \(b/a=\sqrt2-1\) and \(3^{-1/2}\), with \(\nu_2/\nu_1=1.6\) and the line through the centre of the square (numerical). At a fixed toric frame \(T\) is therefore a functional of the edge that depends on its direction relative to the lattice of sites, and it is not the volume of the edge in a quadratic form.

\emph{Rows.} The classicalised charge of~\cite{FlissFrenkelHartnollSoni2025} is \(Q=-2\mathrm i(Y-Y^\dagger)\) with \(Y=\sum_aX^a_{\Sigma\bar\Sigma}\Pi^a_{\bar\Sigma\Sigma}\); it contains no contraction of the directions. The range of \(Y\) lies in the span of the \(r\) edge sites; let \(U\) be the matrix of their basis vectors, \(C=U^\dagger Y\), \(\alpha=CU\) and \(C_\perp=C(1-UU^\dagger)\). For \(2r\le M\) the nonzero spectrum of \(Q\) is that of the Hermitian matrix with blocks \(-2\mathrm i(\alpha-\alpha^\dagger)\), \(-2\mathrm iR^\dagger\), \(2\mathrm iR\), \(0\), where \(R^\dagger R=C_\perp C_\perp^\dagger\); for \(\alpha=0\) it is \(\pm2\) times the singular values of \(C_\perp\) (exact). Take Gaussian momenta, independent in \(a\), with variance \(f(l_1,l_2)\) on the matrices of spins \((l_1,l_2)\) under the two adjoint actions. Then \(\langle YY^\dagger\rangle\) is diagonal on the edge sites, with entry the sum over the links to outside sites \(v\) of \(\kappa\,D_v\), \(D_v=\sum_{j\in\Omega}\langle|\Pi_{vj}|^2\rangle\) (exact); the contraction \(\delta_{IJ}\) in \(E\) is that of this variance. For the cap times the whole second sphere \(D_v=W\equiv\sum_{l_1,l_2}f(l_1,l_2)\,(2l_2+1)N^{-1}w_{l_1}\) on every such link, with \(w_{l_1}\) the sum over the boundary row inside the cap of the squared entries of the spin-\(l_1\) matrices of the first factor, so \(\langle YY^\dagger\rangle=s_1^2W\) on the \(N\) edge sites (exact). At \(N=16\), \(M_1=8\), \(A_f\) multiples of the identity and \(f=\tfrac12l_1\) for \(l_1\le l_{c1}\), \(l_2\le l_{c2}\) and zero otherwise, the sum of the positive eigenvalues of \(Q\), in units of \(2Ns_1W^{1/2}=2W^{1/2}T\), is \(0.85\), \(0.94\), \(0.93\), \(0.95\) for \((l_{c1},l_{c2})=(3,3),(8,3),(15,3),(15,8)\), with sixteen positive eigenvalues whose standard deviation is \(0.3\) of their mean in the last three cases (numerical, \(60\) samples, momenta not constrained to \(\operatorname{Tr}Q=0\)). In this regime the number of boxes of the diagram on each side is therefore close to \(2W^{1/2}T\): it follows \(T\), which here is \(N^{1/2}s\), and not \(s\).

\emph{Dimension.} For the highest weight of \(U(M)\) with \(r\) rows of length \(\ell\) and \(r\) rows of length \(-\ell\), the Weyl dimension formula gives \(\mathrm d\log d/\mathrm d\log M=2r\ell\) to within \(2\%\) for \(M\ge4000\), \(r\le160\), \(\ell\le40\), and \(\log d=2r\ell\,[\log(M/(r+\ell))+c]\) with \(0.80<c<0.97\) for \(1\le r\le160\), \(4\le\ell\le400\), \(M=4000\) and \(16000\); for \(r>1\) this is below \(r\) times the value for one row on each side, by a factor between \(0.69\) and \(0.98\) in that range (numerical). For the cap times the whole second sphere, \(M=M_1N\) and \(r=N\): at \(M_1=50\), \(\ell=30\), \(\log d/(2N\ell)\) rises from \(3.38\) at \(N=10\) to \(4.60\) at \(N=160\), against \(\log M_1=3.91\) (numerical). The rank of the block thus enters through \(M/(r+\ell)\), which for \(\ell\ll N\) is the depth \(M_1\) of the block behind its edge.

\emph{Bearing on Assumption~\ref{ass:Sent-density}.} (i) On this four-dimensional factor the number of boxes, the coefficient of the logarithm in \(\log d_\mu\), is, within the accuracy of the numerical comparison above, a sum over the sites of a three-dimensional edge of a quantity fixed by the links at that site and by the row sums \(D_v\) of the state; the single row of the preceding paragraph was special to a one-dimensional edge. (ii) The measure of that sum is \(T\) weighted by the state. It coincides with a volume, that of \(\tilde\gamma\) up to the factor \(\sqrt2\,(2\pi)^2\), only for an edge along a lattice direction or along the diagonal, with constant speed along the flow; \(\tilde\gamma\) is not the induced tensor, is not \(K^{-1}\), and is not identified with the metric of Equation~\eqref{eq:g-from-e}, whose relation to the matrix configuration rests on Assumption~\ref{ass:dictionary} and Problem~D. (iii) The dependence on the rank remains, through \(\log(M/(r+\ell))\), and the value for \(r\) rows is not \(r\) times the value for one row, so at this order \(\log d_\mu\) is not the integral over the edge of a density fixed by local data of the edge (conditional on the leading-order entropy of~\cite{FlissFrenkelHartnollSoni2025}, at fixed frame). (iv) The row sums \(D_v\), the cut-offs and the frame are supplied, as in the preceding paragraph, and the factor is compact with a positive contraction.

These statements are for the product of two deformed fuzzy spheres with positive contraction, toric blocks at one fixed frame without the excluded pattern, Gaussian momenta with the stated variances, equal rows for the dimension, and the leading semiclassical order of~\cite{FlissFrenkelHartnollSoni2025}; the frame average and its minimisation, projectors that are not toric, a block of a non-compact or Lorentzian factor and the state are not treated. They derive no entropy density on \(\partial M\), no value or scaling of \(S_{\mathrm{ent}}^{\mathrm{(IR)}}\), no boundary stationarity and no \(\Lambda_{\mathrm{eff}}\); Assumptions~\ref{ass:Sent-density} and~\ref{ass:bdy-EE} are not discharged, and Problem~G remains open and PROGRAMME-CONDITIONAL.
